% Шаблон документа для курса "Системы научных публикаций"
% Семинар 11. LaTeX. Создание и редактирование уравнений
\documentclass[a4paper,12pt]{article}

%%% Работа с русским языком
\usepackage{cmap}          % Поиск в PDF
\usepackage{mathtext}         % Русские буквы в фомулах
\usepackage[T2A]{fontenc}      % Кодировка
\usepackage[utf8]{inputenc}      % Кодировка исходного текста
\usepackage[english,russian]{babel}  % Локализация и переносы
\usepackage{indentfirst}            % Для красной строки первого абзаца

\usepackage{extsizes}               % Возможность сделать 14-й шрифт
\usepackage{geometry}               % Простой способ задавать поля
\geometry{top=25mm}
\geometry{bottom=35mm}
\geometry{left=35mm}
\geometry{right=20mm}

%\usepackage{fancyhdr}                   % Колонтитулы
%    \pagestyle{fancy}
%    \renewcommand{\headrulewidth}{0mm}  % Толщина линейки, отчеркивающей
%    \lfoot{Нижний левый}                % верхний колонтитул
%    \rfoot{Нижний правый}
%    \rhead{Верхний правый}
%    \chead{Верхний в центре}
%    \lhead{Верхний левый}
%    % \cfoot{Нижний в центре} % По умолчанию здесь номер страницы (\thepage)


\usepackage{setspace} % Межстрочное расстояние
%\onehalfspacing      % Межстрочное расстояние 1.5
%\doublespacing       % Межстрочное расстояние 2
%\singlespacing       % Межстрочное расстояние 1

\usepackage{lastpage} % Узнать, сколько всего страниц в документе

\usepackage{soul}     % Модификаторы начертания

\usepackage{multicol} % Несколько колонок

\usepackage{lipsum}   % Генерация случайного текста

%\renewcommand{\familydefault}{\sfdefault} % Начертание шрифта без засечек

\usepackage{hyperref}
\usepackage[usenames,dvipsnames,svgnames,table,rgb]{xcolor}
\hypersetup{        % Гиперссылки
  unicode=true,           % Русские буквы в раздела PDF
  pdftitle={Заголовок},   % Заголовок
  pdfauthor={Автор},      % Автор
  pdfsubject={Тема},      % Тема
  pdfcreator={Создатель}, % Создатель
  pdfproducer={Производитель}, % Производитель
  pdfkeywords={keyword1} {key2} {key3}, % Ключевые слова
  colorlinks=true,         % false: ссылки в рамках; true: цветные ссылки
  linkcolor=black,          % Внутренние ссылки
  citecolor=green,        % Ссылки на библиографию
  filecolor=magenta,      % Ссылки на файлы
  urlcolor=blue           % Ссылки на URL
}


%%% Дополнительная работа с математикой
\usepackage{amsfonts,amssymb,amsthm,mathtools} % AMS
\usepackage{amsmath}
\usepackage{icomma} % "Умная" запятая: $0,2$ --- число, $0, 2$ --- перечисление

%% Номера формул
%\mathtoolsset{showonlyrefs=true} % Показывать номера только у тех формул, на которые есть \eqref{} в тексте.

%% Шрифты
\usepackage{euscript} % Шрифт Евклид
\usepackage{mathrsfs} % Красивый матшрифт

%% Перенос знаков в формулах (по Львовскому)
\newcommand*{\hm}[1]{#1\nobreak\discretionary{}
{\hbox{$\mathsurround=0pt #1$}}{}}

\DeclareMathOperator{\sgn}{\mathop{\mathrm{sgn}}}

%%% Окружения типа theorem
\theoremstyle{definition}
\newtheorem{theorem}{Теорема}[section]
\newtheorem{proposition}[theorem]{Утверждение}
\newtheorem{definition}{Определение}[theorem]
\theoremstyle{remark}
\newtheorem*{remark}{Замечание}

%%% Заголовок
\author{Соин Андрей Дмитриевич}
\title{\LaTeX. Создание и редактирование уравнений}
\date{\today}

\begin{document} % Конец преамбулы, начало документа

\maketitle

\section{Формулы и ссылки на них}

Внутритекстовая формула: $a^2 + b^2 = c^2$ — теорема Пифагора.

Выключная формула:
\[
  E = mc^2
\]

Выключная формула с окружением \texttt{equation} и меткой:
\begin{equation}\label{eq:newton}
  F = ma
\end{equation}

Ссылка на формулу: уравнение \eqref{eq:newton} (см.~стр.~\pageref{eq:newton}) — второй закон Ньютона.

\section{Перенос знака в формулах}

Длинный текст длинный текст длинный текст длинный текст
$1\hm{+}2\hm{+}3\hm{+}4\hm{+}5\hm{+}6=21$ длинный текст длинный текст длинный текст.

\section{Формулы}

\begin{equation}
  \int_{-\infty}^{\infty} e^{-x^2} dx = \sqrt{\pi}
\end{equation}

\begin{equation}
  n! \approx \sqrt{2\pi n} \left(\frac{n}{e}\right)^n
\end{equation}

\begin{equation}
  G(t) = B_0 + \sum_{n=1}^{M} \left[ C_n \cos\left(\frac{2\pi n t}{T}\right) + D_n \sin\left(\frac{2\pi n t}{T}\right) \right]
\end{equation}

\begin{equation}
  M\ddot{y} + b\dot{y} + ky = F_1 \cos(2\pi ft)
\end{equation}

\begin{equation}
  \mathrm{CO}_2 + \mathrm{H}_2\mathrm{O} \to \mathrm{H}_2\mathrm{CO}_3
\end{equation}

\begin{multline}
  10+20+30+40+50+60+70+\dots \\
  +500+510+520+530+540+550+560+570+\dots \\
  +960+970+980+990+1000=50500
\end{multline}

\begin{align}
  5 \times 5 &= 25 & \omega_1 + \omega_2 &\notin \Omega \\
  z^2 - 2z + 5 &= 0 & \Gamma \setminus \Lambda &= \Theta \\
  \sqrt{1+y\sqrt{y^3+\sqrt{y+1}}} &= z_0 & \overline{z_1 \cdot z_2} &= \bar{z}_1 \cdot \bar{z}_2
\end{align}

\begin{equation}
  \begin{aligned}
    5 \times 5 &= 25 & \omega_1 + \omega_2 &\notin \Omega \\
    z^2 - 2z + 5 &= 0 & \Gamma \setminus \Lambda &= \Theta \\
    \sqrt{1+y\sqrt{y^3+\sqrt{y+1}}} &= z_0 & \overline{z_1 \cdot z_2} &= \bar{z}_1 \cdot \bar{z}_2
  \end{aligned}
\end{equation}

\begin{equation}
  \begin{aligned}
    I_1 &= 0 \\
    R_1 I_1 + R_2 I_2 + R_3 I_3 &= \mathcal{E}_1 \\
    R_2 I_2 - R_3 I_3 &= \mathcal{E}_2
  \end{aligned}
\end{equation}

\begin{equation}
  A_n =
  \begin{cases}
    1, & \text{если } n=0; \\
    2, & \text{если } n=1; \\
    A_{n-1} - A_{n-2}, & \text{если } n>1.
  \end{cases}
\end{equation}

\begin{equation}
  \exp x =
  \begin{cases}
    e^x, & x>0 \\
    1, & x=0 \\
    e^{-x}, & x<0
  \end{cases}
\end{equation}

\begin{equation}
  \Lambda_x =
  \begin{pmatrix}
    \cosh \chi & -\sinh \chi \\
    -\sinh \chi & \cosh \chi
  \end{pmatrix}
\end{equation}

\begin{equation}
  B =
  \begin{vmatrix}
    b_{11} & \cdots & b_{1m} \\
    \vdots & \ddots & \vdots \\
    b_{m1} & \cdots & b_{mm}
  \end{vmatrix}
\end{equation}

\section{Теоремы}

\begin{theorem}[Теорема косинусов]\label{thm:cos}
  Квадрат стороны треугольника равен сумме квадратов двух других сторон
  без удвоенного произведения этих сторон на косинус угла между ними:
  \[
    c^2 = a^2 + b^2 - 2ab\cos\gamma .
  \]
\end{theorem}

Смотри теорему~\ref{thm:cos} на стр.~\pageref{thm:cos}.

\begin{proposition}\label{prop:twolines}
  Через две точки можно провести две прямые.
\end{proposition}

\begin{definition}
  Биссектрисой угла называется луч, исходящий из вершины угла и делящий его пополам.
\end{definition}

\begin{remark}
  Утверждение~\ref{prop:twolines} не верно.
\end{remark}

\begin{proof}[Доказательство теоремы косинусов]
  \lipsum[1-2]
\end{proof}

\end{document}